% ---- pg-076 （书页 64） ----

如果向量中的一个,如 $\bm{x}_1$ 是 $\bm{0}$,则向量集是线性相关的(令 $a_1$ 是定义 2A.6 中唯一的非零系数).

只有一个元素是 1,而别处是零的向量集是线性无关的.如 $m=4$ 时,

\[
\bm{x}_1 = \begin{bmatrix} 1 \\ 0 \\ 0 \\ 0 \end{bmatrix}, \quad
\bm{x}_2 = \begin{bmatrix} 0 \\ 1 \\ 0 \\ 0 \end{bmatrix}, \quad
\bm{x}_3 = \begin{bmatrix} 0 \\ 0 \\ 1 \\ 0 \end{bmatrix}, \quad
\bm{x}_4 = \begin{bmatrix} 0 \\ 0 \\ 0 \\ 1 \end{bmatrix}
\]

因此,

\[
\bm{0} = a_1\bm{x}_1 + a_2\bm{x}_2 + a_3\bm{x}_3 + a_4\bm{x}_4
= \begin{bmatrix}
a_1 \cdot 1 + a_2 \cdot 0 + a_3 \cdot 0 + a_4 \cdot 0 \\
a_1 \cdot 0 + a_2 \cdot 1 + a_3 \cdot 0 + a_4 \cdot 0 \\
a_1 \cdot 0 + a_2 \cdot 0 + a_3 \cdot 1 + a_4 \cdot 0 \\
a_1 \cdot 0 + a_2 \cdot 0 + a_3 \cdot 0 + a_4 \cdot 1
\end{bmatrix}
= \begin{bmatrix} a_1 \\ a_2 \\ a_3 \\ a_4 \end{bmatrix}
\]

意味着 $a_1 = a_2 = a_3 = a_4 = 0$.

又如,令 $k=3$ 和 $m=3$ 且设

\[
\bm{x}_1 = \begin{bmatrix} 1 \\ 1 \\ 1 \end{bmatrix}, \quad
\bm{x}_2 = \begin{bmatrix} 2 \\ 5 \\ -1 \end{bmatrix}, \quad
\bm{x}_3 = \begin{bmatrix} 0 \\ 1 \\ -1 \end{bmatrix}
\]

则

\[
2\bm{x}_1 - \bm{x}_2 + 3\bm{x}_3 = \bm{0}
\]

因此,$\bm{x}_1 ,\bm{x}_2 ,\bm{x}_3$ 是一个线性相关的向量集,因为其中的任何一个都可以表示成其他向量的线性组合(例如 $\bm{x}_1 = 2\bm{x}_2 + 3\bm{x}_3$).

\mvbold{定义 2A.7} \quad 任意 $m$ 个线性无关的向量集合称为由所有 $m$ 元实数构成的向量空间的一组基.

\mvthm{结论 2A.1}{} \quad 每一个向量可以表示成一组固定基的一个唯一的线性组合.

对于 $m=4$,通常选择的基是

\[
\begin{bmatrix} 1 \\ 0 \\ 0 \\ 0 \end{bmatrix},
\begin{bmatrix} 0 \\ 1 \\ 0 \\ 0 \end{bmatrix},
\begin{bmatrix} 0 \\ 0 \\ 1 \\ 0 \end{bmatrix},
\begin{bmatrix} 0 \\ 0 \\ 0 \\ 1 \end{bmatrix}
\]

已经证明这四个向量是线性无关的.任一向量 $\bm{x}$ 可以唯一地表示为

\[
x_1 \begin{bmatrix} 1 \\ 0 \\ 0 \\ 0 \end{bmatrix} + x_2 \begin{bmatrix} 0 \\ 1 \\ 0 \\ 0 \end{bmatrix} + x_3 \begin{bmatrix} 0 \\ 0 \\ 1 \\ 0 \end{bmatrix} + x_4 \begin{bmatrix} 0 \\ 0 \\ 0 \\ 1 \end{bmatrix} = \begin{bmatrix} x_1 \\ x_2 \\ x_3 \\ x_4 \end{bmatrix} = \bm{x}
\]

由 $m$ 个元素组成的一个向量在几何上可以认为是 $m$ 维空间中的一个点.例如,当 $m=2$,可以认为向量 $\bm{x}$ 是平面上坐标为 $x_1$ 和 $x_2$ 的点.

向量有长度和方向等几何性质.

\mvbold{定义 2A.8} \quad 一个从原点射出的含 $m$ 个元素的向量的长度由勾股定理公式给出:

\begin{mvfigure}
\begin{tikzpicture}[scale=0.9]
\draw[->] (0,0) -- (4.6,0) node[right]{$x_1$};
\draw[->] (0,0) -- (0,2.6) node[left]{$x_2$};
\draw[->] (0,0) -- (3.2,2.0) node[above right]{$\bm{x} = \begin{bmatrix} x_1 \\ x_2 \end{bmatrix}$};
\draw[dashed] (3.2,0) -- (3.2,2.0);
\draw[dashed] (0,2.0) -- (3.2,2.0);
\node at (1.6,-0.35) {$x_1$};
\node at (-0.45,1.0) {$x_2$};
\node[above] at (1.6,1.3) {$2$};
\node[right] at (3.35,1.0) {$1$};
\end{tikzpicture}
\end{mvfigure}
